% ══════════════════════════════════════════════════════════════════════════════
%  EXPERIMENT 7 — FILTRATION                                     (Lesson 15, Term 2 Week 6)
% ══════════════════════════════════════════════════════════════════════════════
\experimentdiv{Filtration}{Folding and using a filter paper \textbullet{} Comparing filters fairly \textbullet{} Residue and filtrate}{Lesson 15 \textbullet\ Term 2, Week 6 \textbullet\ Lab}\label{sec:7}
\lessoncodes{VC2S8U06, U07 \textbullet{} VC2S8I02, I03, I06, I08}

\vocabhead
\begin{ctab}
\begin{tabular}{|SW{4.2cm}|W{11.6cm}|}
  \hline\hd & \thead{Word} & \thead{Definition (copy from the slide)} \\ \hline
  & \scaf{\textbf{Filtration}} & \scaf{Separating an insoluble solid from a liquid by passing it through a filter.} \\ \hline
  & \textbf{Filtrate} & \\ \hline
  & \textbf{Residue} & \\ \hline
  & \textbf{Decanting} & \\ \hline
  & \textbf{Pore} & \\ \hline
\end{tabular}
\end{ctab}

\begin{lipanel}
\inquiryli{Set up filtration correctly and compare three filters in a fair test.}
\understandli{Explain how a filter separates an insoluble solid from a liquid.}
\sctitle
\begin{itemize}[leftmargin=1.4em, itemsep=0pt, topsep=1pt]
  \item I can fold a filter paper and set up a funnel correctly.
  \item I can name the residue and the filtrate.
  \item I can compare filters and explain which works best.
\end{itemize}
\end{lipanel}

\begin{skillsfocus}
\focus{PC2}{Keep the test of each filter fair.}
\focus{PC7}{Set up the filtration equipment correctly.}
\focus{PC6}{Time each filter and rate the filtrate.}
\focus{CO2}{Draw the set-up as a scientific diagram.}
\focus{EV3}{Find sources of error.}
\end{skillsfocus}

% ─── TECHNIQUE ────────────────────────────────────────────────────────────────
\exppart{Technique spotlight}{Blue}

\begin{technique}{Folding a filter paper}{PC7}
\begin{center}
\begin{tikzpicture}[scale=0.62]
  \draw[eq] (0,0) circle (1.2);
  \draw[Ink, dashed, line width=0.7pt] (-1.2,0) -- (1.2,0);
  \node[eqtext, align=center] at (0,-1.9) {1. Fold\\in half};
  \begin{scope}[shift={(4.2,-0.3)}]
    \draw[eq] (-1.2,0) arc[start angle=180,end angle=0,radius=1.2] -- cycle;
    \draw[Ink, dashed, line width=0.7pt] (0,0) -- (0,1.2);
    \node[eqtext, align=center] at (0,-1.6) {2. Fold in\\half again};
  \end{scope}
  \begin{scope}[shift={(8.4,-0.3)}]
    \draw[eq] (0,0) -- (1.2,0) arc[start angle=0,end angle=90,radius=1.2] -- cycle;
    \node[eqtext, align=center] at (0.5,-1.6) {3. Open one layer\\to make a cone};
  \end{scope}
  \begin{scope}[shift={(13.0,-1.1)}]
    \draw[eq] (0,2.4) -- (0.9,0.9) -- (0.9,0) (2.2,2.4) -- (1.3,0.9) -- (1.3,0);
    \draw[Ink, line width=1.1pt, fill=black!10] (0.25,2.25) -- (1.1,1.05) -- (1.95,2.25) -- cycle;
    \node[eqtext, align=center] at (1.1,-0.8) {4. Fit it in the funnel;\\wet it to stick};
  \end{scope}
\end{tikzpicture}
\end{center}
\begin{itemize}[leftmargin=1.4em, itemsep=1pt, topsep=2pt]
  \item Keep the top of the paper \textbf{below} the rim of the funnel, and the liquid \textbf{below} the top of the paper, or the mixture runs round the outside.
  \item Pour down a stirring rod so that the liquid does not splash or tear the paper.
  \item Never poke or stir inside the filter paper.
\end{itemize}
\end{technique}

\begin{workedex}{Residue and filtrate}
Muddy water is filtered. The \textbf{residue} is the soil left on the filter paper. The \textbf{filtrate} is the clearer water that passes through.
\end{workedex}

\begin{wedo}{Residue and filtrate}
Sand and salt water are filtered. Name the residue and the filtrate, and explain why the salt does not stay on the filter paper.
\answerbox{4cm}
\end{wedo}

% ─── PLAN AND SAFETY / METHOD ────────────────────────────────────────────────
\exppart{Plan, safety and method}{Teal}

\textbf{Aim:} To compare how well filter paper, a coffee filter and a cotton cloth filter muddy water.

\task Name the variables. \skr{PC2}
\begin{ctab}
\begin{tabular}{|SW{4.4cm}|W{10.8cm}|}
  \hline\hd & \thead{Variable} & \thead{In this experiment} \\ \hline
  & \scaf{Independent} & \scaf{The type of filter} \\ \hline
  & Dependent (two) & \\ \hline
  & Controlled & \\ \hline
  & Controlled & \\ \hline
\end{tabular}
\end{ctab}

\needspace{8\baselineskip}
\task Predict which filter will give the clearest filtrate, and which will be fastest. Give a reason for each. \skr{QP1}
\answerbox{2.4cm}

\subsubsection*{Equipment}
Filter funnel \textbullet{} retort stand and ring clamp (or a conical flask to sit the funnel in) \textbullet{} 3 beakers \textbullet{} 100\,mL measuring cylinder \textbullet{} stirring rod \textbullet{} filter paper, coffee filter and a square of cotton cloth \textbullet{} 300\,mL of muddy water (stir well before each pour) \textbullet{} stopwatch \textbullet{} white card

\begin{infobox}{Safety}
Wear safety glasses. Glass funnels and beakers break easily: keep them in the middle of the bench. Pour the muddy residue into the bucket, not the sink, so the drain does not block.
\end{infobox}

\subsubsection*{Method}
\begin{enumerate}[itemsep=2pt]
  \item Fold the filter paper and fit it in the funnel. Stand the funnel in a clean beaker.
  \item Stir the muddy water, then measure 50\,mL of it.
  \item Pour it into the filter down the stirring rod and start the stopwatch.
  \item Stop the stopwatch when drips are more than 5 seconds apart. Record the time and the volume of filtrate.
  \item Hold the beaker against the white card and rate the filtrate's clearness from 1 (as muddy as at the start) to 5 (completely clear).
  \item Repeat with the coffee filter and then the cloth, using a clean beaker each time.
\end{enumerate}

% ─── RESULTS ─────────────────────────────────────────────────────────────────
\exppart{Results}{Amber}

\subsubsection*{Table 7.1: Comparing three filters \skr{PC6}}
\begin{ctab}
\begin{tabular}{|SW{3.4cm}|C{2.4cm}|C{2.8cm}|C{2.6cm}|W{3.2cm}|}
  \hline\hd & \thead{Filter} & \thead{Time (s)} & \thead{Volume of filtrate (mL)} & \thead{Clearness (1--5)} & \thead{Residue} \\ \hline
  & \scaf{Paper towel} & \scaf{95} & \scaf{38} & \scaf{3} & \scaf{Thin brown layer} \\ \hline
  & Filter paper & & & & \\ \hline
  & Coffee filter & & & & \\ \hline
  & Cotton cloth & & & & \\ \hline
\end{tabular}
\end{ctab}

\needspace{20\baselineskip}
\task Draw a scientific diagram of your filtration set-up. Label the funnel, filter paper, residue, filtrate and beaker. \skr{CO2}
\drawbox{8cm}{}

% ─── ANALYSIS ────────────────────────────────────────────────────────────────
\exppart{Analysis and evaluation}{Coral}

\task Which filter gave the clearest filtrate, and which was fastest? Use your results.
\answerbox{4cm}

\needspace{12\baselineskip}
\task The clearest filter was not the fastest. Explain why, using the idea of the size of the holes (pores) in the filter.
\answerbox{4.5cm}

\needspace{10\baselineskip}
\task Even the clearest filtrate may not be safe to drink. Explain why. What could be done to make it safe?
\answerbox{4cm}

\needspace{12\baselineskip}
\task Give two sources of error and how each could be reduced. \skr{EV3}
\begin{ctab}
\begin{tabular}{|SW{7.2cm}|W{7.8cm}|}
  \hline\hd & \thead{Source of error} & \thead{How to reduce it} \\ \hline
  & \scaf{The mud settled, so later pours were less muddy} & \scaf{Stir the muddy water for the same time before every pour} \\ \hline
  & & \\ \hline
  & & \\ \hline
\end{tabular}
\end{ctab}

\needspace{10\baselineskip}
\task Decanting is quicker than filtering. Explain when you would decant first and then filter.
\answerbox{4cm}

\begin{yr8box}
\begin{keepext}
\exttask Melbourne's drinking water passes through layers of sand and gravel at a treatment plant. Explain how a sand filter works, and why it has gravel at the bottom and fine sand at the top.
\answerbox{5cm}
\end{keepext}
\end{yr8box}

% ─── SKILL CHECK ─────────────────────────────────────────────────────────────
\exppart{Skill check}{NavyBlue}

\begin{skillcheck}
\scrow{PC2}{I kept the test of each filter fair.}
\scrow{PC7}{I folded the filter paper and set up the funnel correctly.}
\scrow{PC6}{I timed each filter and rated each filtrate.}
\scrow{CO2}{I drew and labelled the filtration set-up.}
\scrow{EV3}{I can name sources of error and how to reduce them.}
\end{skillcheck}

\task Which skill will you focus on in Experiment 8? Explain what you will do differently.
\answerbox{5cm}
